% ===================================================================
%  TABLO No 8 — Hasat. — Av, bağ bozumu, balık tutma.
%  Traduction 2026 d'apres texto/io, controlee sur texto/fr et
%  texto/hi. Consigne : voir texto/tr/10-tablo-01.tex.
%
%  L'appel de note est dans le TITRE de la scene — « Hasat (*) » — et
%  non dans un alinea : c'est ainsi que l'ido le compose.
%
%  LE RENVOI (13) DE LA FAUX MANQUE A L'IDO et le francais le donne.
%  La colonne turque le rend, comme les autres traductions : c'est le
%  deuxieme des trois ecarts declares dans outils/renvoji.py.
% ===================================================================

%%K t08-noto-1 noto
\VUnotes{143.2mm}{%
\textit{(*) Çünkü bu söz belirsiz: keten hasadı mı? üzüm mü? elma mı?
Belki daha iyi olur ki «} \VUgras{mesono} \textit{» kullanılsın, ki} Mondo
\textit{XIV, s. 242'de önerilmiş. Ama şimdiye dek o yeni kökü akademi
benimsememiş ve ido hareketinin genel kurallarına göre tartışılmayı
bekliyor.}%
}

%%K t08-apar-1 apar
\VUsaut{5.0mm}
\VUtitre{15.0pt}{60}{TABLO No 8}
\VUsaut{3.0mm}
\VUcentre{13.2pt}{90}{\textit{Birinci sahne.}}
\VUsaut{3.0mm}
\VUpk{10.2pt}{40}{Hasat (*).}
\VUsaut{4.0mm}
\VUpk{10.2pt}{40}{Kırların görünüşü.}
\VUsaut{3.0mm}

%%K t08-c1-01-1 p
1. --- \VUgras{Yazla} birlikte kırlar bir kez daha görünüşünü değiştirmiş.
Karığa emanet edilen tohum sürdü; topraktan küçük yeşil bir bitki çıktı ve
başağını verdi. Tane oluştu, sonra oldu, ve şimdi hasadı biçmekten başka bir
şey kalmadı.

%%K t08-c1-02-1 p
2. --- \VUgras{Kır çiçekleri} açmış. \VUgras{Peygamber çiçekleri}
\textsuperscript{(1)}, \VUgras{gelincikler} \textsuperscript{(2)} ve
\VUgras{yüksükotları} \textsuperscript{(3)} buğday tarlalarının tekdüze rengine
taze \VUgras{çiçekleriyle} sevinçli bir karşıtlık koyar.

%%K t08-c1-03-1 p
3. --- Meyve ağaçları yaprak vermiş; meyveler olmuş. Küçük
\VUgras{kiraz ağacı} \textsuperscript{(4)} güzel \VUgras{kirazlarla}
\textsuperscript{(5)} yüklü; çocuklar onları koparır ve iştahla yer.

%%K t08-c1-04-1 p
4. --- Sürü artık bütün gününü dışarıda geçirebilir.

%%K t08-c1-04-2 p
\VUgras{Kuzular} \textsuperscript{(6)} büyüyüp güzel \VUgras{koyunlar}
\textsuperscript{(7)} ve güzel \VUgras{koçlar} \textsuperscript{(8)} olmuş,
sık yünlü. Üstüne \VUgras{papatyalar} serpilmiş \VUgras{otlağın}
\textsuperscript{(9)} \VUgras{otunu} rahatça otluyorlar.

%%K t08-c1-04-3 p
Yakında bu hayvanların \VUgras{yünü} kırkılacak, ki ondan giysilerimizin kumaşı
olur.

%%K t08-tit-2 sub
\VUsaut{5.0mm}
\VUpk{10.2pt}{40}{Çoban.}
\VUsaut{3.4mm}

%%K t08-c1-05-1 p
5. --- \VUgras{Çoban} \textsuperscript{(10)} onlara pek dikkat ediyor gibi
benzemez. Tek kaygısı ufuktan geçen fırtına, ve \VUgras{sürü çobanı}
\textsuperscript{(13)}, ki \VUgras{sürahisini} \textsuperscript{(14)}
dudaklarına kaldırmış, ondan da kaygısız görünür.

%%K t08-c1-06-1 p
6. --- Sıcak boğucu; çünkü \VUgras{köpek} \textsuperscript{(15)} de
serinlemeye geliyor; \VUgras{derenin} \textsuperscript{(16)} taze suyunu
yalıyor ve kıyının otları arasına sinmiş zavallı küçük bir \VUgras{kurbağayı}
\textsuperscript{(17)} ürkütüyor. Kurbağa da \VUgras{nilüferlerin}
\textsuperscript{(18)} açtığı o duru suya atlıyor.

%%K t08-c1-07-1 p
7. --- Bir \VUgras{kuyruksallayan} \textsuperscript{(19)} o küçük
\VUgras{kaynağın} \textsuperscript{(20)} yanında yıkanıyor; kaynak iki kayanın
arasından fışkırır ve \VUgras{dikenli çalıların} \textsuperscript{(21)} ile
\VUgras{böğürtlen çalılarının} \textsuperscript{(22)} altına gizlenir.

%%K t08-tit-3 sub
\VUsaut{5.0mm}
\VUpk{10.2pt}{40}{Hasadın biçilmesi.}
\VUsaut{3.4mm}

%%K t08-c1-08-1 p
8. --- Kır çocukları için buğdayın biçilmesi büyük eğlence vakti.
Sevinçli \VUgras{takla} \textsuperscript{(24)} atarlar
\VUgras{saman yığınlarının} \textsuperscript{(23)} üstünde,
\VUgras{biçicilere} \textsuperscript{(26)} yardım ederler, ve onlar
\VUgras{buğday tarlasını} \textsuperscript{(27)} \VUgras{oraklarıyla}
\textsuperscript{(25)} biçerken, ki \VUgras{bileği taşında}
\textsuperscript{(30)} iyice bilenmiş, çocuklar buğdayı devşirip \VUgras{demetler} \textsuperscript{(31)} ya
da \VUgras{küçük bağlar} \textsuperscript{(32)} yaparlar.

%%K t08-c1-09-1 p
9. --- Kimi zaman içlerinden büyüklerine \VUgras{tırpanla}
\textsuperscript{(12)} o \VUgras{altın sapları} \textsuperscript{(28)} biçmeye
izin verilir, ki taneyle dolu ağır \VUgras{başakları} \textsuperscript{(29)}
taşır; küçükler de \VUgras{hayvanları} \textsuperscript{(33)} gözetirken, ki
\VUgras{çayırda} \textsuperscript{(34)} büyük \VUgras{ıhlamurun}
\textsuperscript{(35)} gölgesinde otlar, \VUgras{ağlarıyla}
\textsuperscript{(36)} güzel \VUgras{kelebekler} yakalarlar.

%%K t08-c1-09-2 p
Kimi de \VUgras{arabaya} \textsuperscript{(37)} çıkıp, \VUgras{dirgenle}
\textsuperscript{(38)} kaldırılıp kendilerine verilen demetleri dizerler.

%%K t08-c1-10-1 p
10. --- Bütün \VUgras{ovada} \textsuperscript{(39)}, ki oradan
\VUgras{tarlakuşları} \textsuperscript{(41)} havalanır, büyük bir hareket
var. Herkes hasadın
başında. Ama yoksullar yalnız eski aletleri kullanırken ---
\VUgras{oraklar, tırpanlar} ve \VUgras{çekiç} \textsuperscript{(40)} ---,
varlıklı toprak sahipleri buğdayını ya da \VUgras{çavdarını}
\textsuperscript{(43)} \VUgras{biçerdöverle} \textsuperscript{(42)} biçtirir.
Sonra, demetleri samanlığa taşımadan, orada büyük \VUgras{yığınlar}
\textsuperscript{(44)} kurulur.

%%K t08-c1-11-1 p
11. --- Sonra bir \VUgras{harman makinesi} \textsuperscript{(47)} getirilir, ki
onu bir \VUgras{gezer motor} \textsuperscript{(45)} bir \VUgras{kayış}
\textsuperscript{(46)} aracılığıyla çalıştırır; ve böylece tane
\VUgras{samandan} ayrılır, ekildiği tarlanın kendisinde.

%%K t08-tit-4 sub
\VUsaut{5.0mm}
\VUpk{10.2pt}{40}{Fırtına.}
\VUsaut{3.4mm}

%%K t08-c1-12-1 p
12. --- Hasat vaktinde çok kez sert \VUgras{fırtınalar} \textsuperscript{(53)}
kopar. Önce uzaktan \VUgras{gök gürültüsü} duyulur; sert bir \VUgras{lodos}
\textsuperscript{(54)} kalkar ve \VUgras{ormanın} \textsuperscript{(49)} en
büyük ağaçlarını inleterek eğer. Kara \VUgras{bulutlar}
\textsuperscript{(56)} göğe tırmanır; onları \VUgras{şimşeğin}
\textsuperscript{(55)} göz kamaştıran eğri çizgileri yarar. \VUgras{Yağmur} ve
kimi zaman \VUgras{dolu} düşmeye başlar, ve biçilmeye hazır hasadı yatırır.

%%K t08-c1-13-1 p
13. --- Çok kez yıldırım bir yapıyı tutuşturur. Söz gelişi geçen yıl
\VUgras{yel değirmeninin} \textsuperscript{(50)} çatısı yanıp kül oldu ve iki
\VUgras{kanadı} \textsuperscript{(51)} paramparça oldu.

%%K t08-c1-14-1 p
14. --- Birkaç saat sonra dinginlik döner. Ufukta parlak bir
\VUgras{gökkuşağı} \textsuperscript{(52)} görünür, ve doğa bir süre duran o
yaşamı yeniden alır. \VUgras{Kulübelere} \textsuperscript{(48)} sığınmış olan
köylüler işe döner ve az önceki zararı onarmaya yürekle koyulur.

%%K t08-tit-5 sub
\VUsaut{6.0mm}
\VUcentre{13.2pt}{90}{\textit{İkinci sahne.}}
\VUsaut{3.6mm}

%%K t08-tit-6 sub
\VUpk{10.2pt}{40}{Av. --- Bağ bozumu.}
\VUsaut{1.4mm}
\VUpk{10.2pt}{40}{Balık tutma.}
\VUsaut{4.0mm}

%%K t08-c2-01-1 p
1. --- \VUgras{Ağustos} sonunda ya da \VUgras{eylül} ortasında, yöreye göre, av
mevsimi açılır. Güneşin ilk ışınları \VUgras{kertenkeleleri}
\textsuperscript{(2)} deliklerinden çıkarmaya kalmaz ki \VUgras{avcı}
\textsuperscript{(5)} \VUgras{köpekleriyle} \textsuperscript{(4)} yola çıkar.

%%K t08-c2-02-1 p
2. --- Kalın kumaştan sağlam \VUgras{av takımını} \textsuperscript{(9)}
giymiş, ve deriden \VUgras{tozlukları}, ki \VUgras{karakurbağalarının}
\textsuperscript{(1)} gizlendiği o ıslak otlarda yürümesini sağlayacak;
\VUgras{tüfeğini} \textsuperscript{(6)} ile \VUgras{av çantasını}
\textsuperscript{(10)} almış, ve yürümüş.

%%K t08-c2-03-1 p
3. --- Av hevesi onu bağ bozumunun tam kızıştığı bir bağın ortasına getirir.
Bağcının çocukları, ki çoğu kez yolun kıyısında keçi güderler, bugün onları
canları nerede isterse orada otlatıyorlar, ve yaşlı bir \VUgras{tekeyi}
\textsuperscript{(3)} kazığa bağlayıp --- ki özgürlüğünden yararlanıp elbette
kaçardı --- kendileri de paylarına düşen keyfi sürüyorlar. Biri
\VUgras{üzüm sepetinden} \textsuperscript{(12)} bir salkım alıyor ve
\VUgras{suyunu} \textsuperscript{(14)} bir \VUgras{tasa} \textsuperscript{(13)}
sıkıp şıra (mayalanmamış şarap) yapmakla eğleniyor. Öteki, az önce ördüğü
\VUgras{yaprak tacı} \textsuperscript{(15)} başına koyuyor. Yanlarında arsız
\VUgras{serçeler} \textsuperscript{(16)} üzüm çalmak için cendereye kadar
geliyor.

%%K t08-c2-04-1 p
4. --- Çocuklar eğlenirken adamlar çalışıyor. \VUgras{Üzüm toplayıcıları}
\textsuperscript{(18)} üzümü \VUgras{sırt sepetleriyle} \textsuperscript{(17)}
mahzene taşıyor; bir \VUgras{fıçıcı} \textsuperscript{(19)}
\VUgras{fıçıları} \textsuperscript{(20)} onarıyor, \VUgras{tahtaların}
\textsuperscript{(22)} ve \VUgras{kapakların} \textsuperscript{(21)} sağlam
olduğuna bakıyor, ve kırık \VUgras{çemberleri} \textsuperscript{(23)}
değiştiriyor. O büyük \VUgras{tekneleri} \textsuperscript{(25)} de o yapmıştı, ki
içlerine mayalansın diye \VUgras{üzüm} \textsuperscript{(26)} konur.

%%K t08-c2-05-1 p
5. --- Mayalanma bitince şarap çekilir ve artan \VUgras{cendereye}
\textsuperscript{(28)} konur; büyük \VUgras{vida} \textsuperscript{(29)} ağır
ağır sıkılarak su sıkılır, ki bir \VUgras{tekneye} \textsuperscript{(32)} akar
ve daha da \VUgras{şarap} \textsuperscript{(31)} verir.

%%K t08-tit-8 sub
\VUsaut{6.0mm}
\VUpk{10.2pt}{40}{Bira ve elma şarabı.}
\VUsaut{3.4mm}

%%K t08-c2-06-1 p
6. --- \VUgras{Mahzenin} \textsuperscript{(27)} duvarına bir \VUgras{şerbetçi
otu} \textsuperscript{(30)} dalı tırmanır. Burada yalnız bir süs bitkisi,
ama asmanın pek az yetiştiği kimi ülkelerde \VUgras{bira} yapmak için şerbetçi
otu yetiştirilir.

%%K t08-c2-07-1 p
7. --- Küçük \VUgras{elma ağacı} \textsuperscript{(33)} elmalarla yüklü, ki
olunca güzel \VUgras{elma şarabı} olacak. \VUgras{Kafeslerinde}
\textsuperscript{(34)} \VUgras{kanarya} \textsuperscript{(35)} ile
\VUgras{saka kuşu} \textsuperscript{(36)} özgürce sıçrayıp duran o serçelere
kıskançlıkla bakıyor.

%%K t08-c2-08-1 p
8. --- Göz alabildiğine, tepelerin yamacında, \VUgras{asma kazıklarının}
\textsuperscript{(40)} sıraları dizili, ki \VUgras{salkımlarla} yüklü
görkemli \VUgras{asma kütüklerini} \textsuperscript{(39)} tutar.

%%K t08-c2-08-2 p
Bütün yıl \VUgras{bağcı} \textsuperscript{(42)} \VUgras{bağına}
\textsuperscript{(38)} bakmak için emek vermiş, ve şimdi iyi bir hasatla
ödülünü alır. \VUgras{Toplayıcı kadınlar} \textsuperscript{(41)}
\VUgras{katırların} \textsuperscript{(43)} çektiği arabanın üstünde duran
tekneleri dolduruyor, ve herkes sevinçli.

%%K t08-tit-9 sub
\VUsaut{5.0mm}
\VUpk{10.2pt}{40}{Balık tutma.}
\VUsaut{3.4mm}

%%K t08-c2-09-1 p
9. --- Aynı şey \VUgras{oltacılar} \textsuperscript{(44)} için de doğru; tan
ağarır ağarmaz \VUgras{oltalarıyla} \textsuperscript{(45)} suyun kıyısına
yerleşmişler.

%%K t08-c2-09-2 p
\VUgras{Balıklar} \textsuperscript{(47)} \VUgras{misinalarını}
\textsuperscript{(46)} suya atar atmaz iğneyi yutuyor, ve \VUgras{yemi}
değiştirmeye vakit bulamadan misinanın ucunda yeni bir av çırpınıyor.

%%K t08-c2-09-3 p
Yine de \VUgras{ağcı balıkçı} \textsuperscript{(48)}, ki \VUgras{kayığından}
\textsuperscript{(49)} \VUgras{ırmağın} \textsuperscript{(50)} ortasına ağını
atar, çok daha fazlasını tutar.

%%K t08-c2-10-1 p
10. --- Sonbahar güzel gezinti mevsimi: yaya, \VUgras{at}
\textsuperscript{(52)} üstünde, \VUgras{bisikletle} \textsuperscript{(53)},
\VUgras{otomobille} \textsuperscript{(54)}. \VUgras{Yollar}
\textsuperscript{(51)} temizdir ve insanlar son güzel günlerden tam
yararlanıyor, çünkü \VUgras{kırlangıçlar} \textsuperscript{(56)} sıcak
ülkelere uçmak için şimdiden damlarda toplanmaya başlamış. Yaşlı
\VUgras{ceviz ağacının} \textsuperscript{(57)} yaprakları sararıyor,
\VUgras{cevizler} düşüyor, ve bir \VUgras{koruda} \textsuperscript{(58)}
duran yüksek \VUgras{ağaçlar}, ki \VUgras{yüksek toprakta}
\textsuperscript{(59)} bulunur, yakında çıplak kalacak.

%%K t08-c2-11-1 p
11. --- \VUgras{Güneş} \textsuperscript{(60)} geç doğuyor, günler kısalıyor,
\VUgras{sabahları} iyice bir serinlik duyuluyor, ve \VUgras{yaban ördeği}
\textsuperscript{(61)} sürüleri şimdiden acımasız iklimimizi bırakıp gidiyor.
\VUsaut{6.0mm}
\VUfilet{22mm}
